\section{The transition for every pair of residue classes}
\label{sec:generalthreshold}

Fix distinct $a,b\in G$ and write $B_h(t)=\mathcal P_a(t)-\mathcal P_b(t)$.
Its limiting mean
\[
 M=b_0(a)-b_0(b)-\frac{r_q(a)-r_q(b)}{2d}
\]
is independent of $h$. This includes possible central zeros.
\begin{theorem}\label{th:generaltransition}
If $M=0$, both signs of $B_h$ occur with positive logarithmic density
for every $h>0$. If $M\ne0$, there is a unique $h_*(a,b)>0$ below
which both signs have positive logarithmic density. Above it,
$\operatorname{sgn}(M) B_h$ is bounded below by a positive constant
at all sufficiently late times. At the threshold this oriented race
is positive with logarithmic density one and has lower limit zero.
\end{theorem}

\begin{proof}
Write $f_h$ for the prime-only difference on the phase group $H$,
so $B_h(t)=f_h(\theta_t)+o(1)$ as $t\to\infty$ for fixed $h$.
The semigroup argument of Theorem~\ref{th:transition} shows that
$\min_H f_h$ is strictly increasing and $\max_H f_h$ strictly decreasing:
Lemma~\ref{le:rank} ensures nonconstancy at every $h$.
Both extrema are continuous and converge to $M$ as $h\to\infty$.

To determine the other endpoint, take a prime $p\nmid q$ and put
$t_p=\log p$. We claim
\[
 f_h(\theta_{t_p})=
 \frac{\log p}{\sqrt{4\pi h p}}
  \bigl(\mathbf 1_{p\equiv a\,(q)}-\mathbf 1_{p\equiv b\,(q)}\bigr)
       +O_{p,q}(1) \qquad (h\downarrow0).
\]
This statement concerns the limiting phase function at the fixed time
$t_p$. To prove it, form the tempered distribution
$U_\chi(t)=-\sum_\gamma m_\chi(\gamma)e^{i\gamma t}$
over all signed ordinates, including zero. Temperedness follows from
the zero count. On the open half-line $t>0$, the explicit formula
for the primitive inducer gives the distribution identity
\[
 U_\chi=
 \sum_{n\ge2}\frac{\Lambda(n)\chi^*(n)}{\sqrt n}\,\delta_{\log n}
 -\delta_\chi e^{t/2}
 +\sum_{\omega\ {\rm trivial}}e^{(\omega-1/2)t}.
\]
The last sum runs over the actual trivial zeros of $L(s,\chi^*)$,
with multiplicity, and converges locally to an analytic function.
Indeed test on $\phi\in C_c^\infty(0,\infty)$ and insert
$\int\phi(t)e^{(s-1/2)t}\,dt$ in the logarithmic-derivative contour
of Proposition~\ref{pr:explicit}. Fourier inversion on $\Re s=2$
gives the prime-power atoms. Shift first to $\Re s=-N-1/2$ for an
integer $N$, taking admissible horizontal heights as in that proposition.
The poles and zeros give the terms displayed. On the new line the
functional equation and Stirling's estimate away from trivial zeros
bound $L'/L$ by $O_q(\log(2+N+|\Im s|))$. Integration by parts in the
test transform yields vertical decay of arbitrary fixed order with a
polynomial cost in $N$, times $e^{-N\inf\operatorname{supp}\phi}$.
The vertical remainder therefore vanishes as $N\to\infty$.

Take the character combination for $a-b$ and subtract the constant
prime-square shift. Convolution with $k_h$ gives exactly
$f_h(\theta_t)$ by its absolutely convergent Fourier series. Near
$t_p$ the distribution is a smooth function plus one point mass;
character orthogonality gives the coefficient in the claimed formula,
since $p\nmid q$. Localize by a smooth cutoff equal to one near $t_p$.
The smooth term contributes $O(1)$. The complementary tempered
distribution contributes $o(1)$ because every Schwartz seminorm of the
Gaussian translated to $t_p$ and multiplied by the complementary
cutoff tends to zero exponentially in $1/h$. This proves the claim.

Each of $a$ and $b$ contains a prime. Evaluating at one from each
class proves that the maximum tends to $+\infty$ and the minimum
to $-\infty$ as $h\downarrow0$. If $M\ne0$, orient the race by its
sign and apply strict monotonicity of the minimum. Every nonempty open
sign set has positive Haar mass, while a level set has zero mass by
Theorem~\ref{th:local}. This gives all density and endpoint assertions,
including the zero-mean case by nonconstancy.
\end{proof}

\begin{corollary}[Inversion symmetry]\label{cor:inversion}
If $a\ne a^{-1}$, each class in the prime race $a$ versus $a^{-1}$
leads with logarithmic density exactly $1/2$ for every $h>0$.
In particular this holds for $2$ versus $3$ modulo $5$.
\end{corollary}
\begin{proof}
The character formula gives
$b_\gamma(a^{-1})=\overline{b_\gamma(a)}$ for each signed ordinate,
and $b_0$ is real. Inversion bijects the square roots of $a$ with
those of $a^{-1}$. Thus $z\mapsto\overline z$ exchanges the two
prime-only coordinates on $H$, including their constant terms.
This group involution preserves Haar measure and reverses the
difference. The tie set has measure zero by Theorem~\ref{th:local}.
\end{proof}
